\documentclass[12pt,a4paper]{article}
\usepackage{amsmath, amssymb, amsthm}
\usepackage[utf8]{inputenc}
\usepackage[T1]{fontenc}
\usepackage{geometry}
\usepackage{enumitem}
\usepackage{xcolor}
\geometry{margin=2.5cm}

\newcommand{\easy}{{\color{green!60!black}$\bigstar$}}
\newcommand{\medium}{{\color{orange}$\bigstar\bigstar$}}
\newcommand{\hard}{{\color{red}$\bigstar\bigstar\bigstar$}}

\title{\textbf{Additional Exercises 2}\\
\large General Series --- Absolute \& Conditional Convergence\\[4pt]
\normalsize Mirrors: SIT2 \quad|\quad
\easy\ Easy \quad \medium\ Medium \quad \hard\ Hard}
\date{}

\begin{document}
\maketitle
\thispagestyle{empty}

%----------------------------------------------------------
\section*{Exercise 1: Telescoping series --- prove the identity}
%----------------------------------------------------------

\begin{enumerate}[leftmargin=*, label=(\alph*)]

\item \easy\quad $\displaystyle\sum_{n=1}^{\infty}\frac{1}{n(n+1)} = 1$

\item \easy\quad $\displaystyle\sum_{n=1}^{\infty}\frac{1}{n(n+2)} = \frac{3}{4}$

\item \medium\quad $\displaystyle\sum_{n=2}^{\infty}\ln\!\left(1-\frac{1}{n^2}\right) = -\ln 2$

\item \hard\quad $\displaystyle\sum_{n=1}^{\infty}\frac{1}{n(n+1)(n+2)} = \frac{1}{4}$
\quad\emph{(use partial fractions)}

\end{enumerate}

%----------------------------------------------------------
\section*{Exercise 2: Alternating series --- classify convergence}
%----------------------------------------------------------

\emph{For each series, determine: (i) Does it converge absolutely?
(ii) Does it converge conditionally (Leibniz test)?
(iii) Does it diverge?}

\begin{enumerate}[leftmargin=*, label=(\alph*)]

\item \easy\quad $\displaystyle\sum_{n=1}^{\infty}\frac{(-1)^{n+1}}{n^2+1}$

\item \easy\quad $\displaystyle\sum_{n=1}^{\infty}\frac{(-1)^n}{\sqrt{n}}$

\item \easy\quad $\displaystyle\sum_{n=1}^{\infty}\frac{(-1)^n(n+1)}{n}$

\item \medium\quad $\displaystyle\sum_{n=1}^{\infty}(-1)^n\frac{\ln n}{n}$

\item \medium\quad $\displaystyle\sum_{n=1}^{\infty}\frac{(-1)^n}{n(n+1)}$

\item \medium\quad $\displaystyle\sum_{n=1}^{\infty}(-1)^n\!\left(1-\cos\frac{1}{n}\right)$

\item \medium\quad $\displaystyle\sum_{n=1}^{\infty}(-1)^n\!\left(\sqrt{n+1}-\sqrt{n}\right)$

\item \hard\quad $\displaystyle\sum_{n=1}^{\infty}(-1)^n\!\left(1-\left(\frac{n}{n+1}\right)^n\right)^{1/2}$

\item \hard\quad $\displaystyle\sum_{n=1}^{\infty}\frac{(-1)^n(2n)!}{4^n(n!)^2}$

\end{enumerate}

%----------------------------------------------------------
\section*{Exercise 3: General convergence --- classify each series}
%----------------------------------------------------------

\begin{enumerate}[leftmargin=*, label=(\alph*)]

\item \easy\quad $\displaystyle\sum_{n=1}^{\infty}\frac{\cos(2\pi n)}{n^2}$

\item \easy\quad $\displaystyle\sum_{n=1}^{\infty}\frac{n}{e^n}$

\item \medium\quad $\displaystyle\sum_{n=1}^{\infty}\frac{1}{n\cdot n^a},\quad a>0$

\item \medium\quad $\displaystyle\sum_{n=1}^{\infty}\frac{1}{n\sin(1/n)}$

\item \medium\quad $\displaystyle\sum_{n=1}^{\infty}\left(1+\frac{1}{n}\right)^{n^3}\!\cdot\frac{1}{n^3}$

\item \medium\quad $\displaystyle\sum_{n=1}^{\infty}\frac{(2n)!}{4^n(n!)^2}$

\item \hard\quad $\displaystyle\sum_{n=1}^{\infty}\frac{1\cdot 3\cdot 5\cdots(2n-1)}{2\cdot 4\cdot 6\cdots(2n)}$

\item \hard\quad $\displaystyle\sum_{n=2}^{\infty}\frac{(2n-1)!!}{(2n)!!}\cdot 3^{n/3}$

\end{enumerate}

%----------------------------------------------------------
\section*{Exercise 4: Theoretical problems}
%----------------------------------------------------------

\begin{enumerate}[leftmargin=*, label=(\alph*)]

\item \medium\quad
Given $\sum a_n$ converges absolutely. Prove $\sum a_n/n^2$ converges absolutely.

\item \medium\quad
Given $\sum a_n$ converges. Does $\sum \sin(a_n/n^3)$ necessarily converge absolutely?
Prove or give a counterexample.

\item \hard\quad
Prove the Leibniz test: if $\{a_n\}$ is decreasing, $a_n\ge 0$, and $a_n\to 0$,
then $\sum(-1)^n a_n$ converges, and the error of the $N$-th partial sum is $\le a_{N+1}$.

\item \hard\quad
How many terms of $\displaystyle\sum_{n=1}^{\infty}\frac{(-1)^{n+1}}{n(n+1)^2}$
are needed to approximate the sum with accuracy $10^{-5}$?

\end{enumerate}

\bigskip
\begin{center}
\small\itshape
\easy\ Direct Leibniz or absolute comparison\quad
\medium\ Limit comparison + Leibniz\quad
\hard\ Subtle asymptotic analysis or proof required.
\end{center}

\end{document}
